\documentclass{article}

% Shared preamble — found via Tectonic's --search-path set to the repo root.

% Bibliography (natbib + BibTeX; path relative to this file's directory)
\usepackage{natbib}

\title{Khan Academy: Algebra 1 - Unit 01}
\author{My-MatheMaTeX Study Log}
\date{\today}

\begin{document}

\maketitle
\tableofcontents
\newpage

% -----------------------------------------------------------------------
\section{Overview}

This unit covers fundamentals of evaluating Algebraic expressions from the Khan Academy
Algebra course.

We state definitions precisely and prove the key uniqueness theorem.

% -----------------------------------------------------------------------
\section{Linear Equations in One Variable}

\begin{definition}
  A \textbf{linear equation in one variable} $x$ is an equation of the form
  \[
    ax + b = 0, \qquad a, b \in \mathbb{R},\ a \neq 0.
  \]
\end{definition}

\begin{theorem}[Unique Solution]\label{thm:unique}
  A linear equation $ax + b = 0$ with $a \neq 0$ has exactly one solution,
  namely
  \[
    x = -\frac{b}{a}.
  \]
\end{theorem}

\begin{proof}
  Subtracting $b$ from both sides gives $ax = -b$.
  Since $a \neq 0$, dividing both sides by $a$ is valid and yields
  $x = -b/a$.
  To see uniqueness, suppose $ax + b = 0$ and $ay + b = 0$.
  Subtracting gives $a(x - y) = 0$; since $a \neq 0$ we conclude $x = y$.
\end{proof}

% -----------------------------------------------------------------------
\section{Linear Inequalities}

\begin{definition}
  A \textbf{linear inequality in one variable} $x$ is a relation of the form
  \[
    ax + b \;\square\; 0,
  \]
  where $\square \in \{<,\leq,>,\geq\}$ and $a \neq 0$.
\end{definition}

\begin{note}
  When multiplying or dividing both sides of an inequality by a
  \emph{negative} number, the direction of the inequality reverses.
\end{note}

% -----------------------------------------------------------------------
\section{Worked Exercises}

\begin{exercise}
  Evaluate $xy - y + 3x$. Where $x = 3$ and $y = 2$

  \\
  Substituting into the equation gives

  \[3 \times 2 - 2 + 3(3)\]
  Then
  \[6 - 2 + 9\]
  Working by PEMDAS in which Addition and Subtraction are done in the order they appear left-to-right gives:
  \[4+9 = 13\]
\end{exercise}


\begin{exercise}
  Solve $6a+ 4b - 6$. Where $a = 1$ and $b = 3$
\end{exercise}

\textit{Solution.}
$-2x \geq 6 \implies x \leq -3$ (inequality flips when dividing by $-2$).
The solution set is $(-\infty, -3]$. \qed

% -----------------------------------------------------------------------
\bibliographystyle{plainnat}
\bibliography{../../../bibliography/references}

\end{document}